\section{Model}
\label{sec.model}

We work in the static synchronous setting of~\cite{gkl24} and we keep some of
its notation. This section fixes the notation, states the assumptions, and
restates the facts about typical executions that the analysis uses.

\paragraph{Validators}

There are $n$ parties, called validators, $t$ of which may be Byzantine, that
is, they may arbitrarily deviate from the protocol.

\paragraph{Random Oracle}

We assume access to a random oracle $H \colon \{0,1\}^* \to \{0,1\}^\kappa$,
where $\kappa$ is the security parameter.  Each validator is allowed $q$ queries
to the random oracle $H$ per round. The adversary is allowed $tq$ queries. A
query is successful if its output is below the target $T$. A single query is
successful with probability $p=T/2^\kappa$.

\paragraph{Adversary}

We assume a \emph{Byzantine, rushing} adversary, which may corrupt up to $t$
parties. A rushing adversary sees every message sent by an honest validator
before it is delivered to any other party.

\paragraph{Rounds}

Time proceeds in rounds. Rounds are numbered $1,2,\dots$. For $r \leq s$, we
write $[r,s]$ for the set of rounds $\{r,\dots,s\}$, and $[r,s]=\emptyset$ when
$s<r$.

\paragraph{Network}

We assume a synchronous network among the validators. A message sent by an
honest validator at round $i$ is received by every other validator at the
beginning of round $i+1$. The adversary can see the messages of honest parties
at any point of a round. At the end of each round it fixes the order in which
messages are delivered in the next round; it may reorder messages, spoof their
source or add messages of its own, but it cannot drop a message.

\paragraph{Protocol}

Validators run the Bitcoin backbone protocol of~\cite{gkl24}. Every validator
keeps a local chain, initialized with the genesis block $\gen$. At round $i$ a
validator runs $\fn{maxvalid}$ on its local chain and on the chains it received
at the beginning of round $i$. Function $\fn{maxvalid}$ returns the longest
valid chain among its arguments and breaks ties towards the local chain. The
validator then makes its $q$ queries. On success it appends the new block to the
chain returned by $\fn{maxvalid}$. The resulting chain becomes its local chain
and is sent to all validators. Following~\cite{gkl24}, we say that a validator
\emph{has a chain at round $i$} if this was the local chain it started the round
with. The validator \emph{adopts a chain at round $i$} if this was the return
value of $\fn{maxvalid}$ at round $i$. The chain a validator adopts at round
$i$ is at least as long as the chain it has at round $i$. The chain it has at
round $i+1$ is the chain it adopted at round $i$, extended by at most one block
if the validator was successful at round $i$.

\paragraph{Clients}

Clients do not mine. A client may join at any round and asks for a chain to
adopt. We assume no synchrony between clients and validators. A client does not
know whether it is connected to an honest validator, nor how long it must wait
for an honest message. We model this by letting the adversary choose which
chains a client receives and when.

\paragraph{Chains and the block tree}

A chain is a sequence of blocks, each of which contains the hash of its
predecessor. Every chain begins with the genesis block $\gen$. We write $\ch[i]$
for the $i$-th block of $\ch$, counting from $\ch[0]=\gen$, and we say that
$\ch[i]$ is at height $i$. We write $\ch[i{:}j]$ for the subsequence of $\ch$
at heights $i,\dots,j-1$. We omit $i$ when the range starts at the beginning
of $\ch$, and $j$ when it extends to the end. Negative indices count from the
end. So $\head{\ch}$ is the last block of $\ch$, its \emph{tip}, and
$\prune{\ch}{k}$ is $\ch$ without its $k$ last blocks. When $k\ge\len{\ch}-1$,
$\prune{\ch}{k}$ is the chain consisting of the genesis block alone. We write
$\len{\ch}$ to denote the length of $\ch$, which is the number of blocks of
$\ch$ including the genesis block. The height of $\ch$ is the height of its tip,
$\len{\ch}-1$. We write $\ch_1\preceq\ch_2$ when $\ch_1$ is a prefix of $\ch_2$.
That is, $\ch_1 \preceq \ch_2 \iff \ch_2[{:}\len{\ch_1}] = \ch_1$. Every
validator has a chain of length $1$ at round $1$. A block is honest if it was
computed by an honest validator and adversarial otherwise. We regard the genesis
block as an honest block computed at round $0$.

\begin{definition}[Block tree]\label{def.tree}
    The block tree at round $r$, denoted $\bt_r$, is the set of all blocks
    computed in a round of $[0,r]$. Its edges are the hash references between
    these blocks. A block computed by the adversary belongs to $\bt_r$ whether or
    not it was sent. A chain $\ch$ is in $\bt_r$ if all of its blocks are in
    $\bt_r$.
\end{definition}

No copies occur in a typical execution, which we define below. So every block
of $\bt_r$ has a unique height and $\bt_r$ is a tree.

\begin{remark}\label{rem.membership}
    Every chain held by a client at round $r$ is in $\bt_r$, because the blocks
    of that chain exist at round $r$. Every chain adopted by an honest validator
    at round $r$ is in $\bt_{r-1}$. Indeed, let $B$ be a block of that chain other than the genesis block.
    Either the validator received $B$ at a round at most $r$, and then $B$ was
    sent at round $r-1$ at the latest, or the validator computed $B$ itself. In
    the second case $B$ was computed at a round at most $r-1$, because a
    validator adopts before it queries the oracle.
\end{remark}

\begin{lemma}[Fork choice]\label{lem.fork}
    Let $B$ be an honest block at height $\ell$, computed at round $i$. Then
    some honest validator has a chain of length greater than $\ell$ at round
    $i+1$. Every honest validator adopts a chain of length greater than $\ell$
    at every round $w>i$.
\end{lemma}
\begin{proof}
    If $B$ is the genesis block, then $\ell=i=0$, every validator has a chain of
    length $1$ at round $1$, and every chain has length at least $1$. Otherwise
    the honest validator that computed $B$ appended $B$ at round $i$ to the
    chain it adopted at round $i$, made the result its local chain, and sent it.
    The result has length $\ell+1$, so that validator has a chain of length
    greater than $\ell$ at round $i+1$. Every honest validator receives the
    result at the beginning of round $i+1$, so it is an argument of
    $\fn{maxvalid}$. Hence the chain adopted at round $i+1$ has length greater
    than $\ell$. We claim that the same holds at every round $w>i+1$. The chain a validator has at round $w$ is the chain it adopted at round
    $w-1$, possibly extended by one block. By induction on $w$ it has length
    greater than $\ell$. The chain adopted at round $w$ is at least as long.
\end{proof}

\begin{definition}[Common prefix of two chains]\label{def.meet}
    For chains $\ch_1$ and $\ch_2$, let $\ch_1\meet\ch_2$ be the longest chain
    that is a prefix of both.
\end{definition}

Because a block determines all of its ancestors, $\ch_1$ and $\ch_2$ share a
block at height $i$ if and only if $i<\len{\ch_1\meet\ch_2}$. The quantity
$\len{\ch_1\meet\ch_2}$ is therefore the number of blocks the two chains have in
common.

\begin{definition}[Height schedule]\label{def.schedule}
    A height schedule is a non-dec{\allowbreak}reasing function
    $\hgt\colon\mathbb{N}\to\mathbb{N}$.
\end{definition}

\paragraph{Random variables}

Fix a round $i$. Let $X_i=1$ if some honest validator computes a block at
round $i$, and $X_i=0$ otherwise. Let $Y_i=1$ if exactly one honest validator
computes a block at round $i$, and $Y_i=0$ otherwise. Let $Z_i$ be the number
of successful adversarial queries at round $i$. A round with $X_i=1$ is successful
and a round with $Y_i=1$ is a convergence opportunity. For a set of rounds $S$ we
write $X(S)=\sum_{i\in S}X_i$ and define $Y(S)$ and $Z(S)$ in the same way. We
set $X_0=Y_0=Z_0=0$, so that a set of rounds containing the index $0$ contributes
nothing at that index. The parameter $f=\mathbb{E}[X_i]$ is the probability that
a round is successful.

\paragraph{Assumptions}

We write $\advr=t/(n-t)$ for the ratio of the adversarial to the honest query
budget, and we note that $1+\advr=n/(n-t)$.

\begin{assumption}[Honest majority]\label{asm.majority}
    There exist $\epsilon,\delta,f\in(0,1]$ such that $t\le(1-\delta)(n-t)$ and
    $4f+4\epsilon<\delta$.
\end{assumption}

We call $\delta$ the honest majority advantage.
\Cref{asm.majority} is the honest majority assumption of~\cite{gkl24} with the
condition $3f+3\epsilon<\delta$ replaced by $4f+4\epsilon<\delta$. The
strengthening is used only in \Cref{lem.rho}. Every consequence of the original
assumption that we cite remains available.

\paragraph{Typical executions}

Let $\lambda\ge2/f$ be an integer. An execution is $(\epsilon,\lambda)$-typical
if two conditions hold. First, the bounds of \Cref{lem.typical} hold for every
set $S$ of at least $\lambda$ consecutive rounds. Second, no insertion, copy,
prediction or guess occurs, in the sense of~\cite[Definition 4.3]{gkl24}. An
execution fails to be typical with probability at most
\[
    \typ=4L^2e^{-\Omega(\epsilon^2\lambda f)}+3Q^22^{-\kappa}+[(n-t)L]^22^{-\nu}.
\]
Here $L$ is the total run time of the system, $Q$ bounds the number of oracle
queries, and $\nu$ is the min-entropy of the input contribution function. For
$\lambda=\Omega(\log^2\kappa)$ this probability is negligible in $\kappa$. All
statements below are about typical executions. The following facts are
restated from~\cite{gkl24} without proof.

\begin{lemma}[{\cite[Lemma 4.6]{gkl24}}]\label{lem.typical}
    For every set $S$ of at least $\lambda$ consecutive rounds of a typical
    execution,
    \[
        \begin{gathered}
            (1-\epsilon)f|S| < X(S) < (1+\epsilon)f|S|,
            \qquad Y(S) > (1-\epsilon)(1-f)f|S|,\\
            Z(S) < \left(\frac{\advr}{1-f}+\epsilon\right)f|S| .
        \end{gathered}
    \]
\end{lemma}

\begin{lemma}[{Pairing Lemma, \cite[Lemma 4.1]{gkl24}}]\label{lem.unique}
    Let $B$ be the block at height $j$ of a chain $\ch$, computed by an honest
    validator in a convergence opportunity. Then the block at height $j$ of
    any chain is either $B$ or adversarial.
\end{lemma}

\begin{lemma}[{Chain Growth Lemma, \cite[Lemma 4.2]{gkl24}}]\label{lem.growth}
    If at round $v$ an honest validator has a chain of length $m$, then at every
    round $w\ge v$ every honest validator adopts a chain of length at least
    $m+\sum_{i=v}^{w-1}X_i$.
\end{lemma}

\begin{lemma}[{\cite[Lemma 4.8]{gkl24}}]\label{lem.span}
    In a typical execution, any $k\ge2\lambda f$ consecutive blocks of a chain
    have been computed in more than $k/(2f)$ consecutive rounds.
\end{lemma}
